\section{Experimental Setup, Results, and Evaluation}\label{sec:evaluation}

To assess the effectiveness of the proposed system, we conducted a series of 20 simulation sessions, each representing a distinct synthetic user with a unique persona and product goal. The goal of these experiments is twofold: (i) to verify whether the system can successfully guide users toward target products in natural conversations, and (ii) to evaluate the contribution of each core component. To this end, we designed two ablation scenarios: the first disables the personalized user interaction history, and the second removes the product knowledge graph, thereby eliminating structured path guidance. These comparisons provide insight into the impact of personalization and graph-based reasoning on the overall system performance.

We acknowledge that our evaluation is limited to 20 simulated sessions and lacks validation with real-world users. Although the LLM-based simulation provides a cost-efficient and repeatable approach, it cannot fully capture the complexity of human behavior. In addition, we did not compare against proactive recommendation baselines such as Iterative Preference Guidance (IPG) \cite{bi2024proactive} or multivariate policy learning \cite{jeunen2024multi}. Direct comparison remains difficult because most baselines are designed for offline recommendation or static datasets, whereas our system operates in a dynamic conversational environment that combines multi-turn dialogue with knowledge graph reasoning. Incorporating adapted baselines into future work will help establish a stronger empirical foundation for our contributions.

% We acknowledge that our evaluation is limited to 20 simulated sessions, without validation on real-world user interactions. While the LLM-based simulation offers a cost-efficient and repeatable method, it cannot fully substitute for actual human behavior. Furthermore, we have not compared against existing proactive recommendation baselines such as Iterative Preference Guidance (IPG) \cite{bi2024proactive} or multivariate policy learning \cite{jeunen2024multi}. A direct comparison is challenging because most baseline methods are designed for offline recommendation or static datasets, whereas our system operates in a dynamic conversational environment that integrates multi-turn dialogue and knowledge graph reasoning. Nevertheless, incorporating adapted baselines in future work will provide a stronger empirical position for our contributions.



% Nghiên cứu đã tiến hành thực nghiệm bằng cách chạy 20 phiên mô phỏng, tương ứng với 20 người dùng giả lập có hồ sơ và sản phẩm mục tiêu khác nhau. Để đánh giá vai trò của từng thành phần cốt lõi, chúng tôi thực hiện phương pháp đánh giá theo từng phần (component-wise evaluation) với hai kịch bản chính.

\subsection{Evaluation Metrics}

To evaluate the system's performance from both user-centric and business-driven perspectives, we define the following metrics:

\begin{itemize}
    \item \textbf{History Match Rate:}  
    % The proportion of recommendation turns where the suggested product matches an item from the user's past interactions.
    The proportion of turns in which the system’s suggested question contains a product that appears in the user's past interactions.
    This reflects the system's ability to tailor suggestions based on user history.
    
    \item \textbf{Acceptance Rate:}  
    The proportion of times the user accepts a system-suggested question. This measures user engagement and perceived relevance.
    
    \item \textbf{Answerability Rate:}  
    The proportion of suggested questions that the chatbot is able to answer, indicating the alignment between the Question Guidance Component and the chatbot's knowledge base.
    
    \item \textbf{Goal Completion Rate:}  
    The percentage of sessions that successfully end with the user inquiring about the target product. This is the primary metric reflecting business success.
\end{itemize}

% Để lượng giá sự hữu ích hệ thống theo cả hai hướng cho người dùng và doanh nghiệp, chúng tôi định nghĩa các chỉ số đánh giá như sau.
% \begin{itemize}
%     \item \textbf{P\textsubscript{user\_info} (Tỉ lệ gợi ý liên quan):} Tỉ lệ các lượt hội thoại có gợi ý chứa sản phẩm mà người dùng đã tương tác trong quá khứ. Chỉ số này đo lường mức độ cá nhân hóa.
%     \item \textbf{P\textsubscript{user\_acceptance} (Tỉ lệ chấp nhận gợi ý):} Tỉ lệ người dùng chấp nhận câu hỏi gợi ý của hệ thống.
%     \item \textbf{P\textsubscript{chatbot\_compatibility} (Tỉ lệ tương thích gợi ý):} Tỉ lệ các câu hỏi gợi ý mà chatbot có khả năng trả lời (nằm trong phạm vi tri thức của chatbot).
%     \item \textbf{P\textsubscript{target\_arrival} (Tỉ lệ đạt mục tiêu):} Tỉ lệ các phiên hội thoại kết thúc thành công (người dùng hỏi về sản phẩm mục tiêu). Đây là chỉ số quan trọng nhất để đo lường hiệu quả kinh doanh.
% \end{itemize}

% \subsection{Ablation on Personalization}
\subsection{Impact of Personalization}

% \subsection{Kịch bản 1: Lượng giá vai trò của Cá nhân hóa (User-oriented)}
% In this experiment, we assess the impact of personalization by comparing the full system with a variant that removes user history from the path scoring function. Results in Table~\ref{tab:result_user_oriented} show a dramatic drop in \textit{History Match Rate} from 79.6\% to 2.3\% without access to prior interactions. This indicates that the system becomes almost incapable of recommending products previously shown to interest the user. As a consequence, the \textit{Acceptance Rate} also declines. These results highlight the importance of incorporating user history to generate relevant suggestions and sustain user engagement.


In this scenario, we compare the full system with a variant that removes user interaction history from the path scoring function. Results in Table~\ref{tab:result_user_oriented} show that excluding personalization causes the \textbf{History Match Rate} to plummet from 79.6\% to just 2.3\%, indicating that the system becomes almost entirely unable to recommend products previously aligned with the user’s interests. This reduction also leads to a noticeable drop in the \textbf{Acceptance Rate}, confirming that user history plays a vital role in generating compelling suggestions and maintaining engagement.


% Trong kịch bản này, chúng tôi so sánh hệ thống hoàn chỉnh với một phiên bản hệ thống đã loại bỏ thông tin về lịch sử tương tác của người dùng (\texttt{H}) ra khỏi thuật toán chấm điểm đường đi. Kết quả được thể hiện trong Bảng \ref{tab:result_user_oriented} cho thấy khi loại bỏ thông tin cá nhân hóa, chỉ số \texttt{P\_user\_info} giảm mạnh từ 79.6\% xuống chỉ còn 2.3\%. Điều này chứng tỏ hệ thống gần như không thể gợi ý các sản phẩm mà người dùng đã từng quan tâm. Hệ quả là tỉ lệ chấp nhận gợi ý (\texttt{P\_user\_acceptance}) cũng giảm nhẹ. Điều này khẳng định tầm quan trọng của việc tích hợp lịch sử người dùng để tạo ra các gợi ý hấp dẫn và duy trì sự tương tác.

\begin{table}[!t]
    \centering
    \caption{Comparison with and without user history information.}
    \label{tab:result_user_oriented}

    \begin{tabularx}{\columnwidth}{@{} l 
        >{\centering\arraybackslash}X 
        >{\centering\arraybackslash}X 
        >{\centering\arraybackslash}X 
        >{\centering\arraybackslash}X @{}}
        \toprule
        \textbf{System} & 
        \textbf{HMR} & 
        \textbf{AR} &
        \textbf{ACR} & 
        \textbf{GCR} \\
        \midrule
        Full System & 79.6\% & 94.7\% & 52.0\% & 20.0\% \\
        No User History & 2.3\% & 94.4\% & 46.0\% & 20.0\% \\
        \bottomrule
    \end{tabularx}

    \vspace{2mm}
    \footnotesize
    \textbf{HMR}: History Match Rate \quad
    \textbf{AR}: Answerability Rate \\
    \textbf{AcR}: Acceptance Rate \quad
    \textbf{GCR}: Goal Completion Rate
\end{table}



% \begin{table}[!t]
%     \centering
%     \caption{So sánh hiệu năng khi có và không có thông tin lịch sử người dùng.}
%     \label{tab:result_user_oriented}

%     \begin{tabularx}{\columnwidth}{@{} l >{\centering\arraybackslash}X >{\centering\arraybackslash}X >{\centering\arraybackslash}X @{}}
%         \toprule
%         \textbf{Luồng hệ thống} & 
%         \textbf{$P_{\text{user\_info}}$} & 
%         \textbf{$P_{\text{user\_acceptance}}$} & 
%         \textbf{$P_{\text{target\_arrival}}$} \\
%         \midrule
%         Hệ thống hoàn chỉnh & 0.7963 & 0.52 & 0.2 \\
%         Loại bỏ Lịch sử người dùng & 0.0227 & 0.46 & 0.2 \\
%         \bottomrule
%     \end{tabularx}
% \end{table}




% \subsection{Kịch bản 2: Lượng giá vai trò của Đồ thị tri thức (Business-oriented)}
\subsection{Impact of Knowledge Graph on Business-Oriented Goals}

In this scenario, we remove the knowledge graph component entirely. Instead of guiding the conversation through intermediate entities, the Question Suggestion Agent attempts to create direct recommendations from the user's current query to the target product.

The results in Table~\ref{tab:result_business_oriented} show a substantial decline in system performance when the knowledge graph is removed. In particular, the \textbf{Goal Completion Rate} drops to zero, indicating that the system is unable to guide users to the target product. Without structural support, it lacks logical transitions between products, resulting in abrupt and unconvincing suggestions. For example, the system may suggest a high-end device immediately after the user inquires about a budget option, leading to repeated rejections from the Simulated User Agent.

Interestingly, the \textbf{History Match Rate} increases slightly in this scenario. This is likely because user history is still embedded in the prompt, which may lead the model to overfit to previously mentioned items. However, the improvement is minimal and does not compensate for overall decline in coherence and goal alignment caused by the lack of graph-based reasoning.

The absence of the graph also negatively impacts the \textbf{Answerability Rate}, as the model becomes more prone to hallucinating or producing off-topic suggestions when it is no longer grounded by structured relational context. These findings highlight the pivotal role of the product knowledge graph in generating coherent, goal-aligned recommendations that meet both user expectations and business objectives.


% Furthermore, the Answerability Rate also decreases, highlighting the tendency of the LLM to generate hallucinated or off-topic suggestions when not grounded by the graph. These findings confirm that the knowledge graph is a critical component for achieving coherent recommendations and supporting business objectives.


% Trong kịch bản này, chúng tôi loại bỏ hoàn toàn thành phần đồ thị tri thức. Thay vào đó, Tác tử Gợi ý sẽ cố gắng tạo gợi ý trực tiếp từ sản phẩm người dùng đang hỏi đến sản phẩm mục tiêu mà không có các bước trung gian. Kết quả được thể hiện trong Bảng \ref{tab:result_business_oriented} cho thấy một sự sụt giảm nghiêm trọng trên mọi phương diện. Đáng chú ý nhất, tỉ lệ đạt mục tiêu (\texttt{P\_target\_arrival}) giảm xuống 0\%. Điều này xảy ra vì khi không có đồ thị tri thức, hệ thống không thể tìm ra một con đường logic để kết nối sản phẩm. Các gợi ý trở nên đột ngột và thiếu tự nhiên (ví dụ: đang hỏi về một sản phẩm giá rẻ, lại được gợi ý ngay một sản phẩm cao cấp), dẫn đến việc Tác tử Người dùng liên tục từ chối. Hơn nữa, \texttt{P\_chatbot\_compatibility} cũng giảm, cho thấy LLM có xu hướng "ảo giác" và tạo ra các gợi ý nằm ngoài phạm vi tri thức của chatbot khi không có sự ràng buộc từ đồ thị. Kết quả này chứng minh rằng đồ thị tri thức là thành phần tối quan trọng để đạt được mục tiêu kinh doanh.


\begin{table}[!t]
    \centering
    \caption{Performance comparison with and without the knowledge graph}
    \label{tab:result_business_oriented}
    \begin{tabularx}{\columnwidth}{@{} l 
        >{\centering\arraybackslash}X 
        >{\centering\arraybackslash}X 
        >{\centering\arraybackslash}X 
        >{\centering\arraybackslash}X @{}}
        \toprule
        \textbf{System Variant} & 
        \textbf{HMR} & 
        \textbf{AR} & 
        \textbf{AcR} & 
        \textbf{GCR} \\
        \midrule
        Full System & 79.6\% & 94.7\% & 52.0\% & 20.0\% \\
        Without Knowledge Graph & 81.2\% & 63.3\% & 40.0\% & 0.0\% \\
        \bottomrule
    \end{tabularx}

    \vspace{2mm}
    \footnotesize
    \textbf{HMR}: History Match Rate \quad
    \textbf{AR}: Answerability Rate \\
    \textbf{AcR}: Acceptance Rate \quad
    \textbf{GCR}: Goal Completion Rate
\end{table}


% \begin{table}[!t]
%     \centering
%     \caption{Performance comparison with and without knowledge graph}
%     \label{tab:result_business_oriented}

%     \begin{tabularx}{\columnwidth}{@{} l >{\centering\arraybackslash}X >{\centering\arraybackslash}X >{\centering\arraybackslash}X >{\centering\arraybackslash}X @{}}
%         \toprule
%         \textbf{\shortstack{System\\Variant}} & 
%         \textbf{\shortstack{History\\Match\\Rate}} & 
%         \textbf{\shortstack{Answer-\\ability\\Rate}} &
%         \textbf{\shortstack{Accepta-\\nce Rate}} & 
%         \textbf{\shortstack{Goal\\Completi-\\on Rate}} \\
%         \midrule
%         Full System & 79.6\% & 94.7\% & 52.0\% & 20.0\% \\
%         Without Knowledge Graph & 81.2\% & 63.3\% & 40.0\% & 0.0\% \\
%         \bottomrule
%     \end{tabularx}
% \end{table}


